\section{Conclusion et perspectives}
\label{sec:conclusion2}

\textbf{Synthèse des trois lectures.} La cascade arithmétique des nombres premiers actifs $\{3, 5, 7\}$ de la PT, fixée par le point fixe $\mustar = 15$, admet trois réalisations canoniques de natures complémentaires.
\begin{description}[nosep,leftmargin=*]
\item[Lecture algébro-géométrique (\ref{sec:courbe}).] La cascade est l'unique courbe algébrique dans la classe de Kontsevich--Norbury satisfaisant les conditions (A) d'holonomie polynomiale, (B) de seuil d'activation et (C) d'auto-cohérence. Son genre vaut $\genus(\Cper) = \mustar - 1 = 14$.
\item[Lecture différentielle (\ref{sec:soliton}).] La métrique de Fisher--Bianchi $\gPT$ induite par la cascade satisfait l'équation du soliton de Ricci asymptotique $\Ric + \Hess(f) = \lambda(\mu)\, \gPT$ avec $\lambda(\mu) = -1/\mu^{4} + O(1/\mu^{5})$, coefficient $-1$ exact.
\item[Lecture spectrale (\ref{sec:spectrales}).] L'axiome $\sper = 1/2$ se manifeste dans le résidu de Riemann--von Mangoldt via l'identité triple $1/8 = \sper^{2}/2 = \lambdaH = \DeltaMaslov$ (K2), et la dérivée logarithmique du facteur résiduel admet la décomposition exacte $-d/ds \log R(s) = \sum_p \log p\,[1/(p^{s} - 1) + A_p\,p^{-s}]$ vérifiée à $10^{-25}$ en $\Re(s) > 1$ (A6).
\end{description}

\textbf{Statuts épistémiques différenciés.} Les résultats de cet article se rangent en trois catégories qu'il convient de bien distinguer.
\begin{enumerate}[nosep,leftmargin=*]
\item \emph{Théorèmes prouvés inconditionnellement} : le théorème de classification et la formule du genre (\ref{thm:courbe_classification}, \ref{thm:genre}) ; le théorème du quasi-soliton de Ricci (\ref{thm:soliton_PT}).
\item \emph{Identités prouvées} : K2 (\ref{thm:K2}) est démontrée structurellement comme coïncidence exacte de trois sous-identités indépendantes ; A6 (\ref{thm:A6}) est une identité analytique exacte en $\Re(s) > 1$, vérifiée numériquement à toutes les précisions disponibles.
\item \emph{Programme conjectural} : la conjecture RH-PT (\ref{conj:RH_PT}) est un programme de recherche ouverte, en cours d'instruction~\cite{PT_NewMath_Consolidation}. Sa résolution éventuelle entraînerait l'hypothèse de Riemann.
\end{enumerate}

\textbf{Articulation avec l'article compagnon.} Cet article suppose, pour son intelligibilité complète, la lecture de~\cite{PT_Article1_Foundations}, où sont établis les sept théorèmes structurants $\Tzero$--$\Tsix$ et les trois axiomes de pont prouvés $\BAthree$--$\BAfive$. Les lectures géométrique et spectrale développées ici en sont des prolongements non triviaux : elles ne modifient pas le noyau arithmétique de la PT, mais en révèlent les facettes cohomologique, différentielle et spectrale. Une lecture en sens inverse (de cet article vers l'article 1) est également possible : le genre topologique 14 et l'équation du soliton fournissent des contraintes \emph{a posteriori} sur la cascade qui renforcent l'unicité du choix des premiers actifs $\{3, 5, 7\}$.

\textbf{Perspectives.} Trois directions de recherche immédiates se dégagent de cet article.
\begin{itemize}[nosep,leftmargin=*]
\item \emph{Calcul exhaustif des invariants $\omega_{g,n}^{\Cper}$.} La récursion topologique d'Eynard--Orantin permet, en principe, de calculer tous les invariants $\omega_{g,n}^{\Cper}$ de la courbe de persistance. Les premiers calculs (\cite[\S 8]{PT_GeoFlow_PersistenceCurve}) confirment la substitution structurelle $2\pi^{2}/3 \to (1/2) \sum \gammap{p}$ par rapport à la courbe d'Airy de Mirzakhani. Étendre ce calcul à $\omega_{2,1}^{\Cper}$ et au-delà permettrait de tester de nouvelles identités cascade.
\item \emph{Identité exacte du pont d'échelle $\lambda \leftrightarrow B_0$.} La relation $\lambda(\mu) \sim B_0(\mu)^{-8/3}$ est vérifiée numériquement à $0{,}4\,\%$. L'identité \emph{exacte} (à coefficient près) entre le résidu de Ricci $\lambda$ et le coefficient dominant $B_0$ de la récursion topologique constituerait un pont algébrico-différentiel fort.
\item \emph{Promotion cohomologique de K2.} La dérivation de l'identité K2 à partir d'un invariant cohomologique de l'involution $D : \qplus \leftrightarrow \qminus$ (item K3 du programme~\cite{PT_NewMath_Consolidation}) fournirait une preuve unifiée des trois sous-identités, et un cadre rigoureux pour leur extension à des dimensions différentes de $\mathbb R^{4}$.
\end{itemize}

\textbf{Liens avec les autres applications de la PT.} Les lectures ici présentées sont en dialogue actif avec plusieurs domaines de la phénoménologie PT, en particulier la cosmologie (la métrique de Fisher--Bianchi est l'analogue PT de Bianchi~I, cf.~\cite{companion_M3,companion_M5}), la physique des hautes énergies (le couplage $\lambdaH$ identifié en K2 est l'auto-couplage de Higgs mesurable au LHC, cf.~\cite{companion_M5}), et la théorie analytique des nombres (la conjecture RH-PT relie l'opérateur de survie PT aux zéros de Riemann). Aucune de ces applications n'est développée ici ; elles constituent l'horizon \emph{aval} du présent article.

\textbf{Statut éditorial.} Le couple d'articles~\cite{PT_Article1_Foundations}--présent constitue, à la connaissance de l'auteur, le matériel d'entrée le plus consolidé pour une revue par les pairs externe de la Théorie de la Persistance. Les théorèmes T0--T6, les axiomes de pont BA0--BA5, la dérivation de $\alpha_{\rm EM}$, la courbe de persistance, le quasi-soliton de Ricci et les identités spectrales K2 et A6 y sont énoncés avec leur statut épistémique précis. Les questions restant ouvertes (programme RH-PT, identité exacte $\lambda \leftrightarrow B_0$, promotion cohomologique de K2) sont explicitement marquées comme telles et constituent l'horizon de recherche immédiat du programme.
